\documentclass{article}
\usepackage{amsmath, amssymb, amsthm}

\setlength{\parindent}{0pt}

\newtheorem{claim}{Claim}

\begin{document}

The set $\mathbb{C}$ of complex numbers can be canonically identified with the space $\mathbb{R}^2$ by 
treating each $z = x + iy \in \mathbb{C}$ as a column $(x, y)^T \in \mathbb{R}^2$.

\begin{claim}
    Treating $\mathbb{C}$ as a complex vector space, multiplication by $\alpha = a + ib \in \mathbb{C}$ is 
    a linear transformation in $\mathbb{C}$, and its matrix representation is 
    \[
    \begin{bmatrix}
        \alpha
    \end{bmatrix}
    \]
\end{claim}

\begin{proof}
    Define $M_\alpha(z) = \alpha z$. To show that $M_\alpha$ is a linear transformation, it suffices to 
    verify additivity and homogeneity.

    \medskip

    \textbf{Additivity.} For $z_1, z_2 \in \mathbb{C}$,
    \[
    M_\alpha(z_1 + z_2) = \alpha(z_1 + z_2) = \alpha z_1 + \alpha z_2 = M_\alpha(z_1) + M_\alpha(z_2)
    \]

    \medskip

    \textbf{Homogeneity.} For any scalar $\beta \in \mathbb{C}$ and $z \in \mathbb{C}$,
    \[
    M_\alpha(\beta z) = \alpha(\beta z) = \beta(\alpha z) = \beta M_\alpha(z)
    \]

    \medskip

    Hence $M_\alpha$ is a linear transformation. Since $\dim_\mathbb{C}(\mathbb{C}) = 1$, the standard 
    basis of $\mathbb{C}$ is $\{1\}$. Applying $M_\alpha$ to the basis vector gives
    \[
    M_\alpha(1) = \alpha
    \]
    Therefore, the matrix representation of $M_\alpha$ with respect to the standard basis is 
    \[
    \begin{bmatrix}
        \alpha
    \end{bmatrix}
    \]

    \medskip

    Hence $M_\alpha$ is a linear transformation and its matrix representation is 
    \[
    \begin{bmatrix}
        \alpha
    \end{bmatrix}
    \]
\end{proof}

\begin{claim}
    Treating $\mathbb{C}$ as the real vector space $\mathbb{R}^2$, multiplication by 
    $\alpha = a + ib \in \mathbb{C}$ is a linear transformation in $\mathbb{R}^2$, and its matrix 
    representation is 
    \[
    \begin{bmatrix}
        a & -b \\ 
        b & a
    \end{bmatrix}
    \]
\end{claim}

\begin{proof}
    Define $M_\alpha(z) = \alpha z$. By Claim $1$, $M_\alpha$ is additive and homogeneous for all 
    $\beta \in \mathbb{C}$. Additivity holds regardless of the scalar field, so it transfers immediately. 
    For homogeneity, since $\mathbb{R} \subset \mathbb{C}$, the identity 
    $M_\alpha(\beta z) = \beta M_\alpha(z)$ holds in particular for all $\beta \in \mathbb{R}$. Therefore 
    $M_\alpha$ is real-linear.

    \medskip

    Since $\dim_\mathbb{R}(\mathbb{R}^2) = 2$, the standard basis of $\mathbb{R}^2$ is $\{1, i\}$. Applying 
    $M_\alpha$ to the basis vectors gives
    \[
    M_\alpha(1) = \alpha = a + ib
    \]
    and
    \[
    M_\alpha(i) = i\alpha = ia + i^2b = -b + ia
    \]
    Therefore, the matrix representation of $M_\alpha$ with respect to the standard basis is 
    \[
    \begin{bmatrix}
        a & -b \\
        b & a
    \end{bmatrix}
    \]

    \medskip

    Hence $M_\alpha$ is a linear transformation and its matrix representation is 
    \[
    \begin{bmatrix}
        a & -b \\
        b & a
    \end{bmatrix}
    \]
\end{proof}

\begin{claim}
    Define $T(x + iy) = 2x - y + i(x - 3y)$. Then this transformation is not a linear transformation in 
    the complex vector space $\mathbb{C}$, but if we treat $\mathbb{C}$ as the real vector space 
    $\mathbb{R}^2$, then it is a linear transformation.
\end{claim}

\begin{proof}
    We start by showing that $T$ is real-linear by showing that it is both additive and homogeneous.

    \medskip

    \textbf{Additivity.} For $z_1 = x_1 + iy_1, z_2 = x_2 + iy_2 \in \mathbb{C}$, 
    \[
    \begin{aligned}
        T(z_1 + z_2) &= 2(x_1 + x_2) - (y_1 + y_2) + i((x_1 + x_2) - 3(y_1 + y_2)) \\
        &= 2x_1 + 2x_2 - y_1 - y_2 + i(x_1 + x_2) - 3i(y_1 + y_2) \\
        &= 2x_1 + 2x_2 - y_1 - y_2 + ix_1 + ix_2 - 3iy_1 - 3iy_2 \\
        &= (2x_1 - y_1 + ix_1 - 3iy_1) + (2x_2 - y_2 + ix_2 - 3iy_2) \\
        &= (2x_1 - y_1 + i(x_1 - 3y_1)) + (2x_2 - y_2 + i(x_2 - 3y_2)) \\
        &= T(z_1) + T(z_2)
    \end{aligned}
    \]

    \medskip

    \textbf{Homogeneity.} For any scalar $\beta \in \mathbb{R}$ and $z = x + iy \in \mathbb{C}$, 
    \[
    T(\beta z) = 2\beta x - \beta y + i(\beta x - 3\beta y) = \beta(2x - y + i(x - 3y)) = \beta T(z)
    \]

    \medskip

    Hence $T$ is a linear transformation in the real vector space $\mathbb{R}^2$. Now it remains to show 
    that $T$ is not complex-linear. Let $z = 1$ and $\beta = i$. Then 
    \[
    \begin{aligned}
        T(iz) &= T(i) \\
        &= 2(0) - 1 + i(0 - 3(1)) \\
        &= -1 - 3i \\
    \end{aligned}
    \]
    but
    \[
    \begin{aligned}
        iT(z) &= iT(1) \\
        &= i(2(1) - 0 + i(1 - 3(0))) \\
        &= i(2 + i) \\
        &= -1 + 2i
    \end{aligned}
    \]
    Therefore $T(\beta z) \neq \beta T(z)$ for $\beta = i$, so complex-homogeneity fails. Thus $T$ is not 
    complex-linear.

    \medskip

    Hence $T$ is not a linear transformation in the complex vector space $\mathbb{C}$, but if we treat 
    $\mathbb{C}$ as the real vector space $\mathbb{R}^2$, then it is a linear transformation.
\end{proof}

\end{document}
